% !PRINTED 223
\weberpage{223}
%<<CONT
There are exactly 60 of them. That they are all different
from one another is seen when one represents them in the form:

\weq{11}{\Theta^r = \begin{pmatrix} \varepsilon^r, & 0 \\ 0, & 1 \end{pmatrix},}

\weq{12}{\psi\Theta^r = \begin{pmatrix} 0, & 1 \\ -\varepsilon^r, & 0 \end{pmatrix},}

\weq{13}{\Theta^r\chi\Theta^s = \begin{pmatrix} \varepsilon^r\omega, & \varepsilon^{r-s}\beta \\ 1, & -\varepsilon^{-s}\omega \end{pmatrix},}

\weq{14}{\begin{aligned}
\Theta^r\chi\psi\Theta^s &= \begin{pmatrix} -\varepsilon^{r+s}\beta, & \varepsilon^r\omega \\ \varepsilon^s\omega, & 1 \end{pmatrix} \\
&= \begin{pmatrix} -\varepsilon^r\beta\omega^{-1}, & \varepsilon^{r-s} \\ 1, & \varepsilon^{-s}\omega^{-1} \end{pmatrix},
\end{aligned}}

of which, since $-\omega^2$ is no power of $\varepsilon$, evidently no
two are equal to one another.

The icosahedral group has, as we have seen, 12 five-fold
poles. Each pair of these poles is the pair of poles of four substitutions
of degree 5, which together with the identity form a
five-term cycle. Consequently there are in the group
24 elements of degree 5. Likewise there are 20 three-fold poles,
which lead to 20 substitutions of the third order, and 30 two-fold
poles, which two by two are the poles of one substitution
of the second order each, so that there are 15 substitutions of the second
order. Together with the identity this gives

\[
24 + 20 + 15 + 1 = 60.
\]

The determination of $\beta$ in the substitution $\chi$ now results
from a consideration of the degrees of (13) and (14). To this end
we first form, for an arbitrary linear substitution $S = \begin{pmatrix} a, & b \\ c, & d \end{pmatrix}$,
the second and third repetition:

\[
\begin{aligned}
S^2 &= \begin{pmatrix} a^2 + bc, & b(a+d) \\ c(a+d), & bc + d^2 \end{pmatrix} \\
S^3 &= \begin{pmatrix} a^3 + 2abc + bcd, & b(a^2+bc+ad+d^2) \\ c(a^2+bc+ad+d^2), & d^3+abc+2bcd \end{pmatrix},
\end{aligned}
\]

and if we leave aside the case that $b$ or $c$ vanishes,
which does not come into consideration here, we obtain the
necessary and sufficient condition:
%CONT>>
